\documentclass[10pt, a4paper]{article}

% ── PACKAGES ──────────────────────────────────────────────────
\usepackage[utf8]{inputenc}
\usepackage[T1]{fontenc}
\usepackage{geometry}
\geometry{
  top=2.0cm,
  bottom=2.0cm,
  left=2.2cm,
  right=2.2cm
}
\usepackage{hyperref}
\hypersetup{
  colorlinks=true,
  urlcolor=blue,
  linkcolor=blue,
  citecolor=blue,
  pdftitle={The Dimensional Invariant:
            Sleep Duration, Attractor Basin
            Depth, and Long-Term Survival
            Across the Entire History of Life},
  pdfauthor={Eric Robert Lawson /
             OrganismCore},
  pdfsubject={Reasoning artifact and
              theoretical derivation},
  pdfkeywords={attractor geometry,
               sleep duration,
               sleep topology,
               payment architecture,
               deep-time survival,
               living fossils,
               mass extinction,
               dimensional invariant,
               basin depth,
               Lock strategy,
               Navigator strategy,
               convergent evolution,
               synaptic homeostasis,
               OrganismCore}
}
\usepackage{microtype}
\usepackage{parskip}
\usepackage{booktabs}
\usepackage{array}
\usepackage{tabularx}
\usepackage{longtable}
\usepackage{xcolor}
\usepackage{mdframed}
\usepackage{enumitem}
\usepackage{titlesec}
\usepackage{lmodern}
\usepackage{fancyhdr}
\usepackage{lastpage}
\usepackage{amsmath}
\usepackage{amssymb}
\usepackage{multirow}
\usepackage{tikz}
\usetikzlibrary{arrows.meta, positioning}

% ── COLOURS ───────────────────────────────────────────────────
\definecolor{headergreen}{RGB}{20,100,60}
\definecolor{rulecolor}{RGB}{20,100,60}
\definecolor{claimbox}{RGB}{20,100,60}
\definecolor{notebox}{RGB}{240,248,240}
\definecolor{warningbox}{RGB}{255,248,220}
\definecolor{predbox}{RGB}{235,248,255}
\definecolor{lockbox}{RGB}{240,255,240}
\definecolor{navbox}{RGB}{235,245,255}
\definecolor{unstablebox}{RGB}{255,240,235}
\definecolor{invariantbox}{RGB}{248,240,255}
\definecolor{derivbox}{RGB}{255,252,235}
\definecolor{topologybox}{RGB}{245,240,255}

% ── HEADER / FOOTER ───────────────────────────────────────────
\pagestyle{fancy}
\fancyhf{}
\renewcommand{\headrulewidth}{0.4pt}
\fancyhead[L]{\small\color{headergreen}
  OrganismCore Reasoning Artifact v2.0}
\fancyhead[R]{\small\color{headergreen}
  Dimensional Invariant: Sleep \&
  Survival --- 2026-03-18}
\fancyfoot[C]{\small
  Page \thepage\ of \pageref{LastPage}}

% ── SECTION FORMAT ────────────────────────────────────────────
\titleformat{\section}
  {\large\bfseries\color{headergreen}}
  {\thesection.}{0.5em}{}[
  \vspace{-0.3em}
  \textcolor{rulecolor}
  {\rule{\linewidth}{0.6pt}}]
\titleformat{\subsection}
  {\normalsize\bfseries\color{headergreen}}
  {\thesubsection.}{0.5em}{}
\titleformat{\subsubsection}
  {\small\bfseries}
  {\thesubsubsection.}{0.5em}{}

\setlength{\parskip}{0.4em}
\setlength{\parindent}{0pt}

% ══════════════════════════════════════════════════════════════
\begin{document}

% ── TITLE BLOCK ───────────────────────────────────────────────
\begin{center}
  {\LARGE\bfseries\color{headergreen}
  REASONING ARTIFACT}\\[0.4em]
  {\Large\bfseries
  The Dimensional Invariant:}\\[0.3em]
  {\Large\bfseries
  Sleep Duration, Sleep Topology, Attractor
  Basin Depth, and Long-Term Survival\\
  Across the Entire History of Life}\\[0.5em]
  {\normalsize\itshape
  The oldest surviving lineages sleep
  the least.\\
  The minimum of the sleep spectrum is
  the maximum of the survival spectrum.\\
  The topology of sleep payment is itself
  a geometric prediction of the landscape.\\
  This is not a coincidence. This is the
  geometry.}\\[0.6em]
  \noindent\textcolor{rulecolor}
  {\rule{\linewidth}{1pt}}\\[0.4em]
  {\normalsize
  \textbf{Author:} Eric Robert Lawson\\
  \textbf{Affiliation:} OrganismCore
  (Independent Research)\\
  \textbf{ORCID:}
  \href{https://orcid.org/0009-0002-0414-6544}
  {0009-0002-0414-6544}\\
  \textbf{Contact:}
  \href{mailto:OrganismCore@proton.me}
  {OrganismCore@proton.me}\\
  \textbf{Repository:}
  \href{https://github.com/Eric-Robert-Lawson/
  attractor-oncology}
  {\texttt{github.com/Eric-Robert-Lawson/
  attractor-oncology}}\\
  \textbf{Timestamp:} 2026-03-18\\
  \textbf{Theory paper:}
  \href{https://doi.org/10.5281/zenodo.19094935}
  {\texttt{doi:10.5281/zenodo.19094935}}\\
  \textbf{Document version:} 2.0\\
  \textbf{License:}
  Creative Commons Attribution 4.0
  International (CC BY 4.0)}\\[0.4em]
  \noindent\textcolor{rulecolor}
  {\rule{\linewidth}{1pt}}
\end{center}

\vspace{0.5em}

% ── DERIVATION NOTICE ─────────────────────────────────────────
\begin{mdframed}[
  backgroundcolor=derivbox,
  linecolor=headergreen,
  linewidth=0.8pt,
  innertopmargin=0.6em,
  innerbottommargin=0.6em,
  innerleftmargin=0.8em,
  innerrightmargin=0.8em
]
\textbf{DERIVATION NOTICE}

\smallskip
The hypothesis in this document was derived
from first principles using the attractor
geometry framework ($S + N + G \rightarrow R$)
before literature consultation. The derivation
emerged from the sleep-as-basin-navigation-cost
principle and the sleep sonar instrument
(companion documents, same repository,
2026-03-18), and was extended to a universal
claim about deep-time survival before any
palaeontological or comparative survival
literature was consulted.

\smallskip
This version (v2.0) extends v1.0 by making
explicit a geometric claim that was implicit
in the original derivation but not fully
stated: the \textbf{topology of the sleep
payment process} --- the architecture by which
an organism discharges its basin-occupation
cost --- is itself a geometric prediction of
the framework, not a measurement artefact or
calibration problem. Different landscape
positions do not only produce different sleep
durations. They produce different sleep
\textit{topologies}. This is a stronger and
more complete claim. The crocodilian data
(Kelly \& Lesku, 2015, 2019) makes the
topological dimension empirically explicit
and is recorded here as confirmation of a
predicted geometric property, not as an
anomaly requiring special treatment.

\smallskip
All claims were confirmed against literature
after derivation. The derivation record is
the priority-establishing document.
\end{mdframed}

\vspace{0.8em}

% ── MASTER CLAIM BOX ──────────────────────────────────────────
\begin{mdframed}[
  backgroundcolor=claimbox,
  linecolor=claimbox,
  linewidth=0pt,
  innertopmargin=0.7em,
  innerbottommargin=0.7em,
  innerleftmargin=1.0em,
  innerrightmargin=1.0em
]
\color{white}
\textbf{THE DIMENSIONAL INVARIANT}

\medskip
Across all life, at all scales, across all
of deep time: the total cost of an organism's
attractor basin position must be discharged
through rest. The framework makes two
inseparable predictions:

\medskip
\textbf{Prediction 1 --- Quantity:}
The total cost discharged per cycle
(sleep duration or rest-state equivalent)
tracks the navigational demand of the
organism's landscape position.
Basin depth predicts survival horizon.
The oldest surviving lineages carry the
lowest navigational cost. They discharge
the least.

\medskip
\textbf{Prediction 2 --- Topology:}
The architecture by which cost is
discharged --- the sleep topology --- is
itself a geometric property of the
landscape position, not an independent
biological variable. Different landscape
positions do not merely require different
amounts of sleep. They require different
\textit{kinds} of sleep. The topology is
predicted by the geometry.

\medskip
\textbf{The oldest surviving lineages sleep
the least. How they sleep is the geometry
of where they are.}
\end{mdframed}

\vspace{1em}

% ══════════════════════════════════════════════════════════════
\section{The Framework}
\label{sec:framework}
% ══════════════════════════════════════════════════════════════

\subsection{The Attractor Geometry Framework}

\[
S + N + G \;\rightarrow\; R
\]

\begin{itemize}[itemsep=0.2em]
  \item $S$ = structural substrate (genome
    and encoded build capacity)
  \item $N$ = coherence gradient field
    (external or internal field specifying
    basin occupation)
  \item $G$ = global attractor landscape
    (basin depths, barrier heights,
    accessible states, topology of
    transitions)
  \item $R$ = resulting phenotype or state
\end{itemize}

The framework predicts that an organism's
position in $G$ determines three things
simultaneously:

\begin{enumerate}[itemsep=0.2em]
  \item The total navigational cost of
    occupying that position --- what must
    be discharged.
  \item The survival trajectory --- how
    resistant the basin is to perturbation.
  \item The topology of the discharge
    process --- what architecture of rest
    is geometrically required to pay the
    cost from that position.
\end{enumerate}

These three are not independent. They are
three projections of a single geometric
object: the organism's position in the
landscape $G$.

\subsection{The Core Claim: Cost, Topology,
and Survival as a Single Geometric Object}

The metabolic cost of occupying a position
in an attractor landscape is a function of
the navigational demand imposed by that
position. Navigational demand is determined
by: (1) the number of accessible basins;
(2) the height of barriers between basins;
(3) the frequency of required transitions;
(4) the depth of Basin~$\Omega$ if one
exists; and (5) the within-basin metabolic
or computational intensity of the occupied
basin.

Rest is the generic process by which this
accumulated cost is discharged.
\textbf{The architecture of that rest
process --- its topology --- is not a
separate biological variable. It is a
further projection of the same landscape
geometry.}

An organism in a single deep basin with
zero navigational demand has no synaptic
reconfiguration cost to discharge. Its
rest approaches the minimum and its
topology approaches the simplest possible:
metabolic quiescence or biochemical
circadian rhythm only.

An organism in a symmetric dual-basin
landscape discharging high between-basin
switching cost requires consolidated
bilateral deep sleep to perform the full
synaptic reconfiguration that context
switching demands. The topology must be
bilateral because the cost is bilateral
--- both hemispheres participated in the
navigation and both must be reset.

An organism whose ecological landscape
position is ecologically generalist but
whose metabolic architecture is ectothermic
faces a topological constraint: the
metabolic cost of full bilateral wakefulness
exceeds its energetic budget for most of
the day. The geometric solution is a
distributed intermediate rest state ---
bilateral in EEG signature (because the
navigational cost is moderate and not
requiring hemispheric asymmetry) but
extended across time rather than
concentrated into a single deep-sleep
window. This is not an anomaly. It is the
landscape geometry expressing itself
through the only metabolically feasible
topology available to an ectothermic
Navigator.

The dolphin faces a different but
structurally equivalent constraint: the
cost of full bilateral sleep in an open
ocean environment is mortality. The
landscape geometry resolves this by
dissolving the barrier between waking and
sleep entirely --- the two basins merge
into a permanent intermediate state,
Basin~$\Omega$ of the waking/sleep
landscape itself, maintained via
unihemispheric EEG alternation.

These are not special cases requiring
individual explanation. They are the
framework's topological predictions
working correctly.

\begin{quote}
\itshape
Sleep topology is the geometry of where
you are. Duration is how much the geometry
costs. Survival is how deep the geometry
goes.
\end{quote}

% ══════════════════════════════════════════════════════════════
\section{The Sleep Cost Axes}
\label{sec:axes}
% ══════════════════════════════════════════════════════════════

The total cost discharged through sleep is
decomposable into three axes. Each axis
corresponds to a distinct geometric feature
of the landscape position.

\begin{enumerate}[itemsep=0.5em]

  \item \textbf{Axis~1 --- Between-basin
    switching cost.}
    The organism navigates between multiple
    accessible basins under its own
    motivational field. Each transition
    requires bilateral synaptic
    reconfiguration. Cost accumulates
    continuously during waking.
    \textit{Geometric signature:} High
    behavioural flexibility. Multiple
    accessible basins with low to moderate
    barriers. High sleep, typically
    consolidated bilateral deep sleep.
    \textit{Examples:} Felidae (domestic
    cat navigating wild/domestic dual
    basin), Basin~$\Omega$ organisms.

  \item \textbf{Axis~2 --- Within-basin
    metabolic cost.}
    A single highly expensive basin
    consumes all available metabolic
    resources during waking. Sleep is
    required not for synaptic
    reconfiguration but for metabolic
    resource conservation and recovery.
    \textit{Geometric signature:} Low
    behavioural flexibility. Single basin.
    Complete basin commitment. High sleep
    from resource depletion, not
    navigational cost.
    \textit{Example:} Koala.

  \item \textbf{Axis~3 --- Within-basin
    computational intensity.}
    Extreme neural computation within a
    single highly specialised sensorimotor
    basin during waking. Synaptic
    reconfiguration cost from processing
    intensity, not switching.
    \textit{Geometric signature:} Low
    behavioural flexibility. High
    performance within a single
    specialised basin. High sleep from
    computational cost, not navigational
    switching.
    \textit{Example:} Brown bat
    (echolocation-guided aerial
    interception at millisecond
    resolution for 4 hours requires
    20 hours of synaptic reset).

\end{enumerate}

The Lock organism operates on none of
these axes. Zero navigational demand.
Zero within-basin metabolic surplus
cost. Zero within-basin computational
intensity. The sleep minimum is the
signature of zero cost at every axis
simultaneously.

% ═══════════���══════════════════════════════════════════════════
\section{The Sleep Topologies}
\label{sec:topologies}
% ══════════════════════════════════════════════════════════════

The framework predicts that different
landscape positions produce not just
different sleep durations but different
sleep architectures --- different
topologies of the discharge process.
These topologies are geometric predictions.

\begin{mdframed}[
  backgroundcolor=topologybox,
  linecolor=headergreen,
  linewidth=0.8pt,
  innertopmargin=0.6em,
  innerbottommargin=0.6em,
  innerleftmargin=0.8em,
  innerrightmargin=0.8em
]
\textbf{TOPOLOGY 1: Metabolic quiescence
/ biochemical rhythm only}

\smallskip
\textit{Landscape position:} Deepest
single basin. No nervous system or
minimal neural architecture. Zero
navigational demand. Programme
executes without reconfiguration.

\smallskip
\textit{Discharge process:} No sleep in
any neurological sense. Circadian
biochemical rhythms regulate metabolic
processes only. No synaptic cost to
discharge because no synapses that
reconfigure.

\smallskip
\textit{Duration:} Undefined or
approaching zero.

\smallskip
\textit{Examples:} Sponge, jellyfish
(most species), horseshoe crab.

\smallskip
\textit{Geometric prediction:} This
topology is the signature of the deepest
possible basin. Any organism showing this
topology is at or near the maximum
perturbation-invariance extreme of the
Lock strategy.
\end{mdframed}

\vspace{0.6em}

\begin{mdframed}[
  backgroundcolor=topologybox,
  linecolor=headergreen,
  linewidth=0.8pt,
  innertopmargin=0.6em,
  innerbottommargin=0.6em,
  innerleftmargin=0.8em,
  innerrightmargin=0.8em
]
\textbf{TOPOLOGY 2: Distributed bilateral
rest --- the ectotherm Navigator solution}

\smallskip
\textit{Landscape position:} Moderate
Navigator. Multiple accessible basins with
moderate barriers. Moderate navigational
cost. But ectothermic metabolic
architecture means the energy budget for
full bilateral wakefulness is limited to
a short window per day.

\smallskip
\textit{Discharge process:} True bilateral
deep sleep (EEG-confirmed) for 2--6 hours.
Extended distributed bilateral rest state
(reduced neural activity, bilateral EEG,
one eye open behavioural vigilance
maintained through brainstem-level
alerting, not cortical hemispheric
asymmetry) for 12--16 hours. Total
Navigator-equivalent discharge delivered
across the full inactive period rather
than concentrated into a single window.

\smallskip
\textit{Why bilateral and not
unihemispheric?} Because the navigational
cost is moderate and does not require
hemispheric asymmetry in its discharge.
The switching cost is between ecological
basins, not between waking hemispheres.
Both hemispheres carry the cost equally
and both discharge it simultaneously
across the extended rest period.

\smallskip
\textit{Duration:} 2--6h true sleep plus
12--16h Basin~C rest.
Total equivalent: Navigator range.

\smallskip
\textit{Example:} Crocodilian. Confirmed
bilateral EEG during rest periods (Kelly
\& Lesku, 2015, 2019). One-eye-open
vigilance is behavioural, not
unihemispheric in the neurological sense.
This is the critical distinction.

\smallskip
\textit{Geometric prediction:} This
topology will be found specifically in
ectothermic Navigators --- organisms with
ecological generalism whose metabolic
architecture precludes concentrated
bilateral deep sleep windows. It will not
be found in Lock organisms (too little
cost to distribute) or in endothermic
Navigators (sufficient metabolic budget
for concentrated sleep). The crocodilian
is the confirmed case. Others predicted:
large ectothermic generalist reptiles,
some large ectothermic fish.
\end{mdframed}

\vspace{0.6em}

\begin{mdframed}[
  backgroundcolor=topologybox,
  linecolor=headergreen,
  linewidth=0.8pt,
  innertopmargin=0.6em,
  innerbottommargin=0.6em,
  innerleftmargin=0.8em,
  innerrightmargin=0.8em
]
\textbf{TOPOLOGY 3: Consolidated bilateral
deep sleep}

\smallskip
\textit{Landscape position:} Moderate to
high navigational cost. Endothermic
metabolic architecture with sufficient
budget for full bilateral wakefulness
followed by concentrated sleep window.
Single or dual basin landscape.

\smallskip
\textit{Discharge process:} Concentrated
bilateral deep sleep (NREM and REM)
within a single extended window. Full
hemispheric offline state. Full synaptic
downscaling.

\smallskip
\textit{Duration:} 7--22 hours depending
on axis and intensity.

\smallskip
\textit{Examples:} Rodent, human, koala,
brown bat, lion. The standard mammalian
and avian sleep architecture. Most of the
published comparative sleep literature
describes organisms in this topology.

\smallskip
\textit{Geometric prediction:} This
topology is the signature of endothermic
organisms. Within this topology, duration
discriminates between Axis 1 (switching
cost), Axis 2 (metabolic cost), Axis 3
(computational cost), and the Lock minimum
(zero cost). Duration is a valid direct
proxy for total cost within this topology
only.
\end{mdframed}

\vspace{0.6em}

\begin{mdframed}[
  backgroundcolor=topologybox,
  linecolor=headergreen,
  linewidth=0.8pt,
  innertopmargin=0.6em,
  innerbottommargin=0.6em,
  innerleftmargin=0.8em,
  innerrightmargin=0.8em
]
\textbf{TOPOLOGY 4: Permanent
Basin~$\Omega$ --- dissolved barrier}

\smallskip
\textit{Landscape position:} High
navigational cost in an environment where
full bilateral sleep is lethal. The cost
of the sleep basin's exit barrier has been
inverted --- entering full bilateral sleep
has higher mortality cost than the
accumulated synaptic cost of remaining
partially awake indefinitely.

\smallskip
\textit{Discharge process:} The barrier
between waking and sleep is dissolved.
One hemisphere sleeps while the other
remains active, in continuous alternation.
The organism permanently occupies
Basin~$\Omega$ of the waking/sleep
landscape. Discharge is continuous and
partial rather than periodic and complete.

\smallskip
\textit{Why unihemispheric and not
bilateral?} Because the discharge must
maintain motor and sensory function on
the non-sleeping side at all times. This
is a constraint imposed by the landscape
(open ocean mandatory ventilation plus
multi-directional predation threat), not
by the navigational cost itself. The
navigational cost could in principle be
discharged bilaterally if the environment
permitted it. It cannot. Unihemispheric
is the only geometrically feasible
topology for this landscape position.

\smallskip
\textit{Duration:} Approximately 8 hours
total but never bilateral. Continuous
rotation.

\smallskip
\textit{Examples:} Dolphin, porpoise,
some seals. Some migratory birds during
flight.

\smallskip
\textit{Geometric prediction:} This
topology will be found exclusively in
organisms whose landscape position
combines high navigational cost with a
lethal cost attached to full bilateral
sleep. It requires both conditions
simultaneously. Neither alone produces
the topology.
\end{mdframed}

\vspace{0.6em}

\begin{mdframed}[
  backgroundcolor=topologybox,
  linecolor=headergreen,
  linewidth=0.8pt,
  innertopmargin=0.6em,
  innerbottommargin=0.6em,
  innerleftmargin=0.8em,
  innerrightmargin=0.8em
]
\textbf{TOPOLOGY 5: Fragmented shallow
sleep --- the shallow basin signature}

\smallskip
\textit{Landscape position:} Deep single
basin but with high external pressure
(predation, ecological constraint) making
the sleep basin difficult to occupy for
more than short bursts. The sleep basin
is not shallow due to low cost --- it is
shallow because the waking basin is
maintained under high external pressure
that continually pulls the system back
before deep sleep can complete.

\smallskip
\textit{Discharge process:} Brief
fragmented bilateral sleep bouts.
Incomplete discharge per bout. Discharge
accumulated across many short sessions.

\smallskip
\textit{Duration:} 30 minutes to 2 hours
total, distributed across the day.

\smallskip
\textit{Examples:} Giraffe. Migratory
birds during peak migration (Basin~A so
deep that sleep is deferred entirely until
migration terminates, then discharged as
sleep rebound).

\smallskip
\textit{Geometric prediction:} Sleep
rebound is the diagnostic signature of
this topology. An organism in Topology~5
that is temporarily relieved of the
external pressure will show immediate
and disproportionate sleep increase ---
paying accumulated deferred debt. This
distinguishes it from a Lock organism
(which will not show sleep rebound under
safety conditions because it has no
accumulated debt) and from an endothermic
Navigator (which discharges continuously
and has no large debt to rebound from).
\end{mdframed}

\vspace{0.8em}

The five topologies are not a taxonomy of
sleep types. They are the geometric
predictions of the framework applied to
the waking/sleep sub-landscape within $G$.
Each topology is the only architecturally
feasible solution for its corresponding
landscape position. The framework does not
say organisms ``choose'' a sleep topology.
It says the landscape position geometrically
determines which topology is accessible.

% ══════════════════════════════════════════════════════════════
\section{The Two Survival Strategies}
\label{sec:strategies}
% ══════════════════════════════════════════════════════════════

The deep-time survival data reveals two
geometrically opposite strategies that
both produce long-term lineage persistence.

\subsection{Strategy~1: The Lock}

\begin{mdframed}[
  backgroundcolor=lockbox,
  linecolor=headergreen,
  linewidth=0.8pt,
  innertopmargin=0.5em,
  innerbottommargin=0.5em,
  innerleftmargin=0.8em,
  innerrightmargin=0.8em
]
\small
\begin{tabular}{p{3.2cm}p{9.3cm}}
\textbf{Basin structure} &
Single basin. Extremely deep.
Very high exit barrier.\\[0.2em]
\textbf{Morphology} &
Extreme conservatism. Body plan
unchanged across hundreds of millions
of years.\\[0.2em]
\textbf{Metabolism} &
Slow. Efficient. Minimal resource
requirement.\\[0.2em]
\textbf{Behaviour} &
Low flexibility. Narrow repertoire.
Highly programmatic.\\[0.2em]
\textbf{Sleep quantity} &
Minimal. Approaching zero in simplest
forms. 2--5 hours in complex Lock
organisms.\\[0.2em]
\textbf{Sleep topology} &
Topology~1 (quiescence/biochemical
rhythm) in simplest forms.
Topology~5 (fragmented, predation
pressure) in complex Lock organisms
with environmental threat.
Never Topology~4 (no open-ocean
lethal constraint).
Never Topology~2 (insufficient
navigational cost to distribute).\\[0.2em]
\textbf{Survival horizon} &
Maximum. 300--600\,Myr+.\\[0.2em]
\textbf{Survival mechanism} &
Basin so deep it is perturbation-invariant.
Extinction events cannot displace the
organism.\\[0.2em]
\textbf{Convergence} &
Multiple independent phyla.\\
\end{tabular}
\end{mdframed}

\subsection{Strategy~2: The Navigator}

\begin{mdframed}[
  backgroundcolor=navbox,
  linecolor=headergreen,
  linewidth=0.8pt,
  innertopmargin=0.5em,
  innerbottommargin=0.5em,
  innerleftmargin=0.8em,
  innerrightmargin=0.8em
]
\small
\begin{tabular}{p{3.2cm}p{9.3cm}}
\textbf{Basin structure} &
Multiple basins. Shallow barriers.
High switching capacity.\\[0.2em]
\textbf{Morphology} &
Ecological generalism. Body plan
adapted for flexibility.\\[0.2em]
\textbf{Metabolism} &
Flexible. Moderate. Endothermic
or ectothermic --- but the
metabolic architecture determines
which sleep topology is available
to discharge the Navigator cost.\\[0.2em]
\textbf{Behaviour} &
High flexibility. Broad repertoire.
Contextual switching.\\[0.2em]
\textbf{Sleep quantity} &
Moderate equivalent. The total
Navigator-range discharge cost
is constant across Navigators
but is delivered through different
topologies depending on metabolic
architecture.\\[0.2em]
\textbf{Sleep topology} &
Endothermic Navigator:
Topology~3 (consolidated bilateral
deep sleep, 8--14h).
Ectothermic Navigator:
Topology~2 (distributed bilateral
rest, 2--6h true sleep $+$ 12--16h
Basin~C rest, total equivalent
Navigator range).
Open-ocean Navigator with lethal
bilateral constraint:
Topology~4 (permanent Basin~$\Omega$,
unihemispheric alternation).\\[0.2em]
\textbf{Survival horizon} &
Long. 100--300\,Myr.\\[0.2em]
\textbf{Survival mechanism} &
Navigates to alternative basin in
real time when current basin disappears.
Generalism is itself Basin~$\Omega$ of
ecological flexibility.\\[0.2em]
\textbf{Convergence} &
Multiple independent lineages.\\
\end{tabular}
\end{mdframed}

\subsection{The Unstable Zone}

\begin{mdframed}[
  backgroundcolor=unstablebox,
  linecolor=red!40!black,
  linewidth=0.8pt,
  innertopmargin=0.5em,
  innerbottommargin=0.5em,
  innerleftmargin=0.8em,
  innerrightmargin=0.8em
]
\small
\begin{tabular}{p{3.2cm}p{9.3cm}}
\textbf{Basin structure} &
Shallow single basin. High barriers
to alternative basins.\\[0.2em]
\textbf{Morphology} &
Narrow specialist. Not deep enough
to be perturbation-invariant. Not
flexible enough to reach alternative
basins when niche disappears.\\[0.2em]
\textbf{Sleep quantity} &
Variable. Often high (Axis~2 or 3)
but with no survival advantage.\\[0.2em]
\textbf{Sleep topology} &
Typically Topology~3 (consolidated
bilateral) but the duration reflects
niche-maintenance cost, not
navigational capacity. The topology
does not discriminate Unstable Zone
organisms from Navigator organisms
in duration alone --- the
discriminator is behavioural
flexibility and dietary breadth,
not sleep architecture.\\[0.2em]
\textbf{Survival horizon} &
Short. Millions of years.\\[0.2em]
\textbf{Examples} &
\textit{T.~rex}, mammoth, sabre-tooth
cat, passenger pigeon, dodo.\\
\end{tabular}
\end{mdframed}

\subsection{Why The Lock Outlasts The Navigator}

The Lock survives longer than the Navigator
for a precise geometric reason. The Navigator
requires accessible alternative basins at
the moment of disruption. Catastrophic
perturbations that eliminate all accessible
basins simultaneously defeat the Navigator.
The Lock cannot be displaced because its
basin depth exceeds the perturbation energy
of any event the biosphere has yet produced.

The Permian--Triassic event killed 96\% of
all marine species. The horseshoe crab was
still there on the other side.

\begin{quote}
\itshape
Basin depth is the ultimate survival
parameter. Navigation is a survival
strategy for moderate perturbations.
For catastrophic perturbations, only the
deepest basins survive.
The minimum of the sleep spectrum is
the maximum of the survival spectrum.
The topology of that minimum is the
geometry of a basin so deep it needs
no maintenance.
\end{quote}

% ══════════════════════════════════════════════════════════════
\section{The Empirical Record}
\label{sec:empirical}
% ══════════════════════════════════════════════════════════════

\subsection{Lock Strategy: Confirmed Examples}

\begin{center}
\small
\begin{longtable}{p{2.2cm}p{1.4cm}p{3.2cm}p{3.0cm}p{2.8cm}}
\toprule
\textbf{Lineage} &
\textbf{Age} &
\textbf{Sleep quantity} &
\textbf{Sleep topology} &
\textbf{Basin geometry}\\
\midrule
Sponge
(Porifera) &
$\sim$600\,Myr &
Undefined. No nervous
system. &
Topology~1.
Circadian
biochemical only. &
Deepest possible.
Sessile. Single
programme.\\[0.4em]
Jellyfish
(Cnidaria) &
500--700\,Myr &
Undefined for most.
\textit{Cassiopeia}
rest $\sim$8h. &
Topology~1 for
most. Topology~5
fragment in
\textit{Cassiopeia}. &
No brain. Radial
symmetry. Passive
drift and sting.\\[0.4em]
Nautilus
(Mollusca) &
$\sim$500\,Myr &
No REM. Quiescent
periods distributed. &
Topology~1 to
Topology~2 border.
Bilateral rest.
No hemispheric
architecture. &
Neutral buoyancy.
Passive ambush.
Shell unchanged.\\[0.4em]
Horseshoe crab
(Chelicerata) &
$\sim$450\,Myr &
Tidal/circadian
rest only. &
Topology~1.
No bilateral
sleep. &
Deepest terrestrial
Lock. 450\,Myr
unchanged.\\[0.4em]
Shark
(Chondrichthyes) &
$\sim$400\,Myr &
No true bilateral
sleep. Rest periods
$\sim$2--4h
distributed. &
Topology~2 border.
Distributed rest.
No REM confirmed. &
Cartilage. Generalist
enough to border
Lock/Navigator
by species.\\[0.4em]
Coelacanth
(Sarcopterygii) &
$\sim$400\,Myr &
Unknown. Presumed
minimal. &
Topology~1 to~2.
Deep cold water
suppresses all
activity. &
Minimal metabolism.
Unchanged 400\,Myr.\\[0.4em]
Tuatara
(Lepidosauria) &
$\sim$200\,Myr &
12h+ rest. Most
is torpor-like
low-temp dormancy. &
Topology~2.
Distributed
ectotherm rest.
Not true bilateral
deep sleep. &
Slowest reptile
metabolism.
Functions at
5\textdegree C.
Lives 100+ years.\\
\bottomrule
\end{longtable}
\end{center}

\subsection{Navigator Strategy: Confirmed Examples}

\begin{center}
\small
\begin{longtable}{p{2.2cm}p{1.4cm}p{3.2cm}p{3.0cm}p{2.8cm}}
\toprule
\textbf{Lineage} &
\textbf{Age} &
\textbf{Sleep quantity} &
\textbf{Sleep topology} &
\textbf{Basin geometry}\\
\midrule
Cockroach
(Blattodea) &
$\sim$300\,Myr &
$\sim$12h rest
cycle. &
Topology~2.
Distributed
bilateral rest.
Ectotherm. &
Eats anything.
Lives anywhere.
300\,Myr survival.\\[0.4em]
Crocodilian
(Archosauria) &
$\sim$200\,Myr &
2--6h true bilateral
EEG sleep. 12--16h
Basin~C distributed
bilateral rest.
Total equivalent:
Navigator range. &
\textbf{Topology~2.}
Bilateral EEG
throughout (Kelly
\& Lesku, 2015,
2019). One-eye-open
vigilance is
\textit{behavioural},
not unihemispheric
in the neurological
sense. The predicted
topology for an
ectothermic Navigator. &
Generalist predator.
Survived K--Pg.
Fresh/salt water.
Any prey. 200\,Myr.\\[0.4em]
Dolphin
(Cetacea) &
$\sim$50\,Myr &
$\sim$8h total.
Never bilateral. &
\textbf{Topology~4.}
Permanent
Basin~$\Omega$.
Unihemispheric
alternation.
Waking/sleep
barrier dissolved. &
Open ocean.
High cognitive
complexity. Must
surface to breathe.
Bilateral sleep
lethal.\\[0.4em]
Rodent
(Muridae) &
$\sim$65\,Myr &
$\sim$13h. &
Topology~3.
Consolidated
bilateral deep
sleep. Standard
endothermic
Navigator. &
Extreme generalist.
All biomes and
diets.\\[0.4em]
Corvidae &
$\sim$30\,Myr &
$\sim$8--10h. &
Topology~3.
Standard avian
bilateral. Some
Topology~4 during
roosting. &
Problem solving.
Tool use. Social
complexity.\\
\bottomrule
\end{longtable}
\end{center}

\subsection{Unstable Zone: Confirmed Examples}

\begin{center}
\small
\begin{tabular}{p{2.4cm}p{1.4cm}p{3.0cm}p{3.0cm}p{2.8cm}}
\toprule
\textbf{Lineage} &
\textbf{Age} &
\textbf{Sleep quantity} &
\textbf{Sleep topology} &
\textbf{Why extinct}\\
\midrule
\textit{T.~rex}
(Tyranno\-sauridae) &
2--3\,Myr &
Unknown. Moderate
large predator
presumed. &
Topology~3
presumed. &
Narrow apex
predator. No
alternative basin
at K--Pg.\\[0.3em]
Mammoth
(Proboscidea) &
0.5--5\,Myr &
Likely $\sim$2h
(cf. elephant). &
Topology~5
(size + predation
pressure). &
Grassland specialist.
Habitat contracted.
Could not adapt.\\[0.3em]
Sabre-tooth cat
(Machairo\-dontinae) &
0.5--11\,Myr &
High felid-range
likely. &
Topology~3.
Basin~$\Omega$
felid ancestry. &
Specialist predator.
Megafauna gone.
Basin collapsed.\\[0.3em]
Dodo
(Raphidae) &
0.003\,Myr &
Standard avian
$\sim$8--10h. &
Topology~3. &
Island specialist.
Single shallow
basin. First
perturbation
$=$ extinction.\\
\bottomrule
\end{tabular}
\end{center}

% ══════════════════════════════════════════════════════════════
\section{The Sleep Spectrum as Temporal Sonar}
\label{sec:sonar}
% ═════════════════════════════════════════════��════════════════

\subsection{The Full Sonar Instrument}

The sleep sonar has two channels, not one.
Duration alone is insufficient. Duration
and topology together triangulate the
organism's position in the landscape with
full geometric precision.

\begin{center}
\small
\begin{tabular}{p{2.0cm}p{2.0cm}p{2.5cm}p{4.5cm}p{2.0cm}}
\toprule
\textbf{Duration} &
\textbf{Topology} &
\textbf{Survival horizon} &
\textbf{Landscape position} &
\textbf{Examples}\\
\midrule
$\sim$0 &
T1 &
600\,Myr+ &
Deepest Lock. No
nervous system. Zero
cost, zero discharge. &
Sponge, jellyfish\\[0.3em]
0--4h &
T1--T2 &
300--500\,Myr &
Deep Lock. Minimal
neural architecture.
Programme runs
without maintenance. &
Horseshoe crab,
shark, coelacanth\\[0.3em]
2--6h true
$+$ T2 rest &
T2 &
100--300\,Myr &
Ectothermic Navigator.
Navigator cost
distributed across
extended bilateral
rest. Total equivalent:
Navigator range. &
Crocodilian,
cockroach,
tuatara\\[0.3em]
8--14h
concentrated &
T3 &
10--100\,Myr &
Endothermic Navigator
or coupled social
field. Standard
bilateral discharge. &
Rodent, corvid,
dog, human\\[0.3em]
8h, never
bilateral &
T4 &
50--200\,Myr &
Open-ocean Navigator.
Lethal bilateral
constraint. Basin~$\Omega$
of waking/sleep
landscape. &
Dolphin\\[0.3em]
0.5--2h
fragmented &
T5 &
Variable --- external
pressure dependent &
Single deep basin
under strong external
pressure. Sleep
basin shallow by
external force,
not low cost.
Sleep rebound
diagnostic. &
Giraffe\\[0.3em]
14--22h &
T3 &
Niche-dependent &
Axis 2 (metabolic),
Axis 3 (computational),
or deep Basin~$\Omega$
(Felidae). &
Koala, bat,
lion, cat\\
\bottomrule
\end{tabular}
\end{center}

\subsection{The U-Shaped Survival--Duration
Relationship}

The survival--sleep duration relationship is
U-shaped with a flat minimum. The minimum
of the sleep duration spectrum corresponds
to the maximum of the survival spectrum
(Lock strategy). The moderate range is the
Navigator survival zone. The high-duration
range is the uncertain zone --- survival
depends on niche stability or
Basin~$\Omega$ depth.

\textbf{Critical note on duration alone:}
Duration without topology is ambiguous.
A crocodilian sleeping 2--6 hours of true
bilateral sleep appears to be in the Lock
range. It is not. It is an ectothermic
Navigator using Topology~2 to distribute
Navigator-range cost. The topology
disambiguates.

Similarly, a giraffe sleeping 30 minutes
appears to be a deep Lock. It may be.
But the diagnostic is sleep rebound under
safe conditions. If the giraffe increases
sleep substantially when predation threat
is removed, it is a Topology~5 organism
with accumulated debt, not a Lock. If it
does not increase sleep under safety,
the basin depth is geometric and the
Lock classification holds.

\textbf{Duration $+$ Topology $+$ Rebound
response: these three together uniquely
specify the organism's landscape position.
None of the three alone is sufficient.}

% ══════════════════════════════════════════════════════════════
\section{Convergence Across Independent Phyla}
\label{sec:convergence}
% ══════════════════════════════════════════════════════════════

\subsection{The Lock Converges Independently}

\begin{center}
\small
\begin{tabular}{p{3.5cm}p{3.0cm}p{2.0cm}p{3.5cm}}
\toprule
\textbf{Phylum / Class} &
\textbf{Example} &
\textbf{Age} &
\textbf{Sleep topology}\\
\midrule
Porifera & Sponge & $\sim$600\,Myr & T1\\
Cnidaria & Jellyfish & 500--700\,Myr & T1\\
Mollusca & Nautilus & $\sim$500\,Myr & T1--T2 border\\
Chelicerata & Horseshoe crab & $\sim$450\,Myr & T1\\
Chondrichthyes & Shark & $\sim$400\,Myr & T2 border\\
Sarcopterygii & Coelacanth & $\sim$400\,Myr & T1--T2\\
Lepidosauria & Tuatara & $\sim$200\,Myr & T2\\
\bottomrule
\end{tabular}
\end{center}

Seven independent phyla. Seven independent
arrivals at the same geometric strategy.
The Lock topology is Topology~1 at the
extreme and transitions toward Topology~2
as neural complexity increases while
remaining in a single deep basin. This
transition within the Lock is itself a
geometric prediction: as Lock organisms
develop more complex nervous systems, they
acquire more synaptic cost to discharge,
but the basin remains deep enough to keep
total cost minimal. Topology shifts from
T1 to T2 but duration remains in the
Lock range.

\subsection{The Navigator Converges Independently,
Across Multiple Topologies}

\begin{center}
\small
\begin{tabular}{p{3.2cm}p{2.5cm}p{2.0cm}p{4.3cm}}
\toprule
\textbf{Lineage} &
\textbf{Example} &
\textbf{Age} &
\textbf{Sleep topology}\\
\midrule
Blattodea & Cockroach & $\sim$300\,Myr &
T2 (ectotherm)\\
Archosauria & Crocodilian & $\sim$200\,Myr &
T2 (ectotherm, confirmed bilateral EEG)\\
Cephalopoda & Octopus & $\sim$300\,Myr &
T2--T3 border (distinct active/quiet
sleep phases confirmed)\\
Cetacea & Dolphin & $\sim$50\,Myr &
T4 (open ocean lethal bilateral constraint)\\
Rodentia & Rat & $\sim$65\,Myr &
T3 (endotherm)\\
Aves & Corvid & $\sim$30\,Myr &
T3, T4 during roosting\\
Primates & Human & $\sim$6\,Myr &
T3 (endotherm, inter-generational
Basin~B* coupling)\\
\bottomrule
\end{tabular}
\end{center}

The critical geometric observation: the
Navigator strategy has converged
independently across at least three
distinct sleep topologies (T2, T3, T4).
The strategy is the same --- ecological
generalism as Basin~$\Omega$ of ecological
flexibility. The topology varies with the
metabolic and environmental constraints
of the specific lineage. This is exactly
what the framework predicts: the topology
is determined by the sub-constraints of
the landscape position, not by the
strategy classification.

The Navigator strategy is the deep basin.
The topology is the solution the landscape
provides for paying the cost of occupying
that basin given the full constraint set.

\subsection{The Unstable Zone Does Not Converge}

Mass extinction casualties do not converge.
Each extinction event eliminates different
taxa via different mechanisms. The unstable
zone has no deep basin by definition. The
only convergent feature of the unstable
zone is its instability. Topology does not
discriminate here: Unstable Zone organisms
occupy Topology~3 (same as endothermic
Navigators) but lack the dietary breadth
and behavioural flexibility that makes
the Navigator strategy viable. Duration
and topology together still fail to
distinguish Unstable Zone from Navigator
without the behavioural flexibility
measurement --- which is why the full
three-channel sonar (duration $+$ topology
$+$ flexibility) is the complete instrument.

% ══════════════════════════════════════════════════════════════
\section{The Human Position}
\label{sec:human}
% ══════════════════════════════════════════════════════════════

The human is an endothermic Navigator
with Topology~3 sleep (consolidated
bilateral deep sleep, 7--9 hours) and
an unprecedented inter-generational
Basin~B* --- a co-constructed social
attractor that persists beyond the
individual lifetime via cultural
transmission.

The human sleep topology (T3) is standard
endothermic Navigator. The sleep duration
(7--9 hours) is standard Navigator range.
Nothing in the human's sleep signature
is anomalous with respect to the
framework.

What is anomalous is the Basin~B* depth
increasing across generations through
cultural accumulation. This creates a
new geometric feature not present in any
other known organism: a basin whose depth
is a function of time and investment
rather than purely of evolutionary
selection. The framework does not yet
predict the topological consequences
of a basin that deepens in real time
at inter-generational scale. This is
the honest limit of the current derivation.

The risk: the same cultural Basin~B*
is remodelling the ecological landscape
faster than any prior organism. A
Navigator survives by finding alternative
basins when the current one disappears.
If the Navigator is simultaneously the
agent collapsing the alternative basins
for all other species, the Navigator
strategy fails even for the Navigator.
This is a geometric statement, not a
moral one.

% ══════════════════════════════════════════════════════════════
\section{Literature Confirmations}
\label{sec:confirmations}
% ══════════════════════════════════════════════════════════════

\subsection{Oldest Surviving Lineages Sleep
the Least}
\textit{Derived before literature:}
Minimum cost $=$ minimum discharge $=$
maximum basin depth $=$ maximum survival.

\textit{Confirmed:} Horseshoe crab,
nautilus, jellyfish, shark. All Lock
organisms confirmed as minimal discharge.
All show Topology~1 or Topology~2.
Source: Comparative sleep literature;
living fossil biology.

\subsection{Crocodilian Sleep is Bilateral
EEG, Not Unihemispheric}
\textit{Derived before literature:}
Ectothermic Navigator topology prediction
--- distributed bilateral rest, not
unihemispheric.

\textit{Confirmed:} Kelly \& Lesku (2015)
EEG study of juvenile \textit{Crocodylus
niloticus}: bilateral EEG throughout rest
periods. One-eye-open vigilance is
behavioural, maintained via brainstem
alerting, not cortical hemispheric
asymmetry. This is the predicted T2
topology for an ectothermic Navigator.
Source: Kelly, M.L. \& Lesku, J.A. (2015),
\textit{J. Exp. Biol.}, 218, 3175--3180;
Kelly \& Lesku (2019), \textit{Biol. Rev.}

\subsection{Mass Extinction Survivors Are
Disproportionately Lock or Navigator}
\textit{Derived:} Unstable zone is where
extinction concentrates casualties.

\textit{Confirmed:} K--Pg survivors:
crocodilians (Navigator, T2), turtles
(partial Lock, T2), birds (Navigator, T3
and T4), small mammals (Navigator, T3).
Permian--Triassic survivors: horseshoe
crab (Lock, T1), nautilus (Lock, T1).
Source: Jablonski (1986);
Barnosky et al. (2011).

\subsection{Both Strategies Convergent
Across Independent Phyla}
\textit{Derived:} Lock and Navigator are
deep basins in evolutionary strategy
space, therefore convergent.

\textit{Confirmed:} Convergence confirmed
in phyla listed in Section~\ref{sec:convergence}.
Source: Conway Morris (2003,
\textit{Life's Solution}).

\subsection{Slow Metabolism Correlates With
Deep-Time Survival}
\textit{Derived:} Lock strategy requires
minimal metabolic expenditure.

\textit{Confirmed:} Tuatara, coelacanth,
horseshoe crab. All show Lock topology
(T1--T2) and slow metabolism as convergent
traits.
Source: Comparative metabolism literature;
Cree (2014) tuatara biology.

\subsection{Generalist Diet and Ecological
Flexibility Predict Extinction Survival}
\textit{Derived:} Navigator mechanism is
generalism as Basin~$\Omega$.

\textit{Confirmed:} Capellini et al. (2008,
\textit{Proc Royal Soc B}); Barnosky et al.
(2011).

% ══════════════════════════════════════════════════════════════
\section{The Novel Predictions}
\label{sec:predictions}
% ══════════════════════════════════════════════════════════════

All predictions stated before data
gathered to test them. Locked here.
2026-03-18.

\subsection{P-INVARIANT-1: Sleep Duration
Negatively Correlates With Lineage Age
Across All Phyla, Controlling For Topology}

A cross-phyla analysis of sleep duration
(or rest-state equivalent, topology-
calibrated) against lineage age will show
a negative correlation surviving correction
for body mass, metabolic rate, and
topology class. The oldest lineages will
have the lowest topology-calibrated
discharge requirement. Detectable across
Porifera, Cnidaria, Mollusca, Arthropoda,
and Vertebrata simultaneously.

\textit{Discriminating value:} This
prediction is not generated by any
existing sleep function theory. It
requires both the dimensional invariant
and the topology-calibration framework.

\subsection{P-INVARIANT-2: Mass Extinction
Survivor Lineages Sleep Less (Lock) or Show
Higher Dietary Breadth (Navigator) Than
Casualty Sister Taxa}

For each major mass extinction event,
Lock survivors sleep less than casualties.
Navigator survivors show higher dietary
breadth than casualties and occupy T2,
T3, or T4 (not T5 under chronic
ecological pressure). Testable via fossil
record body plan conservatism proxy,
dietary isotope and dental analysis,
and phylogenetic comparison.

\subsection{P-INVARIANT-3: The Unstable Zone
Boundary Is Predictable From Duration,
Topology, and Dietary Breadth}

A three-dimensional boundary in the space
of (sleep duration, sleep topology,
dietary breadth) separates Lock, Navigator,
and Unstable Zone organisms. The two-
dimensional boundary (duration $\times$
dietary breadth) of v1.0 is confirmed
but the addition of topology as the third
axis removes the ambiguity between
ectothermic Navigators (T2, appearing
Lock-range in duration alone) and true
Lock organisms.

\subsection{P-INVARIANT-4: Topology~2 Will
Be Found Specifically and Exclusively in
Ectothermic Navigators}

Topology~2 (distributed bilateral rest
with no EEG hemispheric asymmetry) will
be found in organisms satisfying both:
(a) ecological generalism (Navigator
strategy), and (b) ectothermic metabolic
architecture. It will not be found in Lock
organisms (insufficient navigational cost
to distribute) or in endothermic Navigators
(sufficient metabolic budget for
concentrated sleep). The crocodilian is
the confirmed case. Predicted additionally:
large ectothermic generalist reptiles (e.g.,
monitor lizards, large iguana species),
and large ectothermic generalist fish.

\subsection{P-INVARIANT-5: Sleep Rebound
Uniquely Identifies Topology~5 And
Distinguishes It From Lock}

Organisms showing minimal sleep
(30 minutes to 2 hours) will divide
into two classes under safe-condition
testing:
\textit{Class A (Lock):} No significant
sleep increase when predation pressure
is removed. No rebound. Basin is
geometrically deep. Minimal sleep is the
geometric minimum, not a deferred debt.
\textit{Class B (Topology~5):} Significant
sleep increase when pressure is removed.
Rebound present. Sleep was suppressed by
external pressure, not absent by geometric
necessity.
The giraffe is the predicted Class~B
case. Migratory birds post-migration are
the confirmed Class~B case.

\subsection{P-INVARIANT-6: Lock Organisms
Are Undomesticable}

The same basin depth that produces
perturbation-invariance in the face of
mass extinctions renders the basin
immovable by any domestication field
that does not also destroy the organism.
No Lock organism (T1 or T2 topology,
minimal duration, Lock-range dietary
breadth) has been domesticated. This
will remain true.

\subsection{P-INVARIANT-7: The Dimensional
Invariant Extends Below Nervous System
to Metabolic Rest State Duration in
Single-Celled Organisms}

Single-celled organisms show metabolic
rest states whose duration tracks niche
breadth in the same direction the
invariant predicts. Specialist bacteria
(narrow niche, deep metabolic basin)
show longer continuous activity with
minimal rest. Generalist bacteria show
moderate distributed rest. This extends
the invariant below the level of neural
architecture entirely --- confirming
it as a property of basin topology,
not of sleep as a neurological phenomenon.

% ══════════════════════════════════════════════════════════════
\section{The Dimensional Invariant:
Complete Formal Statement}
\label{sec:statement}
% ══════════════════════════════════════════════════════════════

\begin{mdframed}[
  backgroundcolor=invariantbox,
  linecolor=headergreen,
  linewidth=1pt,
  innertopmargin=0.8em,
  innerbottommargin=0.8em,
  innerleftmargin=1.0em,
  innerrightmargin=1.0em
]
\textbf{THE DIMENSIONAL INVARIANT:
COMPLETE STATEMENT}

\medskip
Across all life, at all scales, across
all of deep time, the framework
$S + N + G \rightarrow R$ makes two
inseparable predictions about rest:

\medskip
\textbf{1. Quantity:} The total cost of
basin occupation (decomposed into Axis~1
between-basin switching cost, Axis~2
within-basin metabolic cost, and Axis~3
within-basin computational intensity)
must be discharged per cycle. Duration
of discharge tracks total cost. Basin
depth --- resistance to perturbation ---
is the primary predictor of survival.
The oldest surviving lineages have the
lowest basin-occupation cost. They
discharge the least.

\medskip
\textbf{2. Topology:} The architecture
of the discharge process --- the sleep
topology --- is a further geometric
prediction of the same framework. The
landscape position determines not only
how much must be discharged but what
architecture of discharge is geometrically
feasible given the full constraint set
(metabolic architecture, environmental
lethality of full sleep, navigational
demand structure).

\medskip
Five topologies are predicted:
T1 (quiescence/biochemical rhythm) for
the deepest Lock organisms;
T2 (distributed bilateral rest) for
ectothermic Navigators;
T3 (consolidated bilateral deep sleep)
for endothermic organisms at all cost
levels;
T4 (permanent Basin~$\Omega$,
unihemispheric) for organisms where
full bilateral sleep is lethal;
T5 (fragmented shallow sleep with rebound)
for organisms under chronic external
suppression of sleep.

\medskip
The invariant holds because basin depth
is the same geometric property at every
scale. Sleep --- in whatever topology the
landscape geometry permits --- is the
payment. Survival is the depth.

\medskip
\textbf{The oldest surviving lineages
sleep the least.}\\
\textbf{How they sleep is the geometry
of where they are.}\\
\textbf{This is not a coincidence.}\\
\textbf{This is the geometry.}
\end{mdframed}

% ══════════════════════════════════════════════════════════════
\section{The Unified Derivation Chain}
\label{sec:chain}
% ══════════════════════════════════════════════════════════════

\begin{center}
\small
\begin{tabular}{p{1.4cm}p{5.5cm}p{5.5cm}}
\toprule
\textbf{Step} &
\textbf{Claim} &
\textbf{Record}\\
\midrule
1. Cancer &
FOXA1/EZH2 ratio identifies whether
a cancer cell is epigenetically arrested
in the cancer attractor or has lost
access to the identity basin. &
\href{https://doi.org/10.5281/zenodo.18883922}
{\texttt{10.5281/zenodo.18883922}}\\[0.4em]
2. Domestication &
The domestication syndrome is
field-specified developmental arrest
of NCC migration. The pig is the
clean case. &
\href{https://doi.org/10.5281/zenodo.19094935}
{\texttt{10.5281/zenodo.19094935}}\\[0.4em]
3. Canine &
Above a genomic fixation threshold,
the landscape is remodelled. Basin~A
becomes inaccessible. &
Repository 2026-03-18\\[0.4em]
4. Feline &
The symmetric dual-basin landscape.
The three-class domestication taxonomy.
Basin~$\Omega$ as meta-attractor. &
Repository 2026-03-18\\[0.4em]
5. Sleep cost &
Sleep duration tracks the total cost
of basin occupation. Three axes.
Universal principle. &
Repository 2026-03-18\\[0.4em]
6. Sleep sonar &
The sleep spectrum maps onto the
three-axis framework across all animals. &
Repository 2026-03-18\\[0.4em]
7. This document &
The dimensional invariant. Sleep quantity
and sleep topology are both geometric
predictions of the landscape. The oldest
surviving lineages sleep the least.
The topology of how they sleep is the
geometry of where they are. &
This document\\
\bottomrule
\end{tabular}
\end{center}

\bigskip

\begin{mdframed}[
  backgroundcolor=claimbox,
  linecolor=claimbox,
  linewidth=0pt,
  innertopmargin=0.7em,
  innerbottommargin=0.7em,
  innerleftmargin=1.0em,
  innerrightmargin=1.0em
]
\color{white}
\textbf{THE UNIFIED STATEMENT}

\medskip
The attractor geometry framework is a
single geometric principle operating
identically at every scale at which
life organises itself.

\medskip
From a cancer cell deciding whether to
differentiate or proliferate,
to a domestic pig deciding whether to
feral,
to a cat choosing which basin to occupy
this morning,
to a crocodile distributed across
16 hours of bilateral rest because its
metabolic budget cannot afford the
concentrated sleep that its navigational
cost demands,
to a dolphin that has dissolved the
barrier between waking and sleep and
now permanently inhabits the intermediate,
to a horseshoe crab that has been in
the same basin for 450 million years
and does not sleep at all because there
is nothing to reset:

\medskip
the geometry is the same.\\
The basin is the same concept.\\
The cost of navigation is the same cost.\\
The sleep is the same payment.\\
The topology of sleep is the shape of
the landscape.\\
The survival is the depth.

\medskip
Scale is the only variable.\\
The geometry is invariant.
\end{mdframed}

% ── LINKED RECORDS ────────────────────────────────────────────
\vspace{0.8em}
\noindent\textcolor{rulecolor}
{\rule{\linewidth}{0.6pt}}

\subsection*{Linked Records}

\begin{center}
\small
\begin{tabular}{p{3.5cm}p{9.0cm}}
\toprule
\textbf{Record} & \textbf{DOI / Location}\\
\midrule
Theory paper &
\href{https://doi.org/10.5281/zenodo.19094935}
{\texttt{doi:10.5281/zenodo.19094935}}\\[0.2em]
P7 pre-registration &
\href{https://doi.org/10.5281/zenodo.19096128}
{\texttt{doi:10.5281/zenodo.19096128}}\\[0.2em]
P9 pre-registration &
{[DOI pending --- insert after
confirmation]}\\[0.2em]
FOXA1/EZH2 cancer &
\href{https://doi.org/10.5281/zenodo.18883922}
{\texttt{doi:10.5281/zenodo.18883922}}\\[0.2em]
Plant inversion &
\href{https://doi.org/10.5281/zenodo.19040399}
{\texttt{doi:10.5281/zenodo.19040399}}\\[0.2em]
LECA resurrection &
\href{https://doi.org/10.5281/zenodo.18986790}
{\texttt{doi:10.5281/zenodo.18986790}}\\[0.2em]
Canine / feline /
sleep artifacts &
\texttt{github.com/Eric-Robert-Lawson/
attractor-oncology} (2026-03-18)\\
\bottomrule
\end{tabular}
\end{center}

% ── DOCUMENT METADATA ─────────────────────────────────────────
\vspace{0.5em}
\noindent\textcolor{rulecolor}
{\rule{\linewidth}{0.8pt}}
\vspace{0.3em}

{\footnotesize
\textbf{Document ID:}
DIMENSIONAL\_INVARIANT\_SLEEP\_SURVIVAL\_
REASONING\_ARTIFACT\_v2.0\\
\textbf{Date:} 2026-03-18\\
\textbf{Author:} Eric Robert Lawson /
OrganismCore\\
\textbf{Status:} REASONING ARTIFACT ---
DERIVATION RECORD.
Version 2.0 extends v1.0 by making
explicit the topology dimension as a
geometric prediction of the framework,
not a measurement correction.
Primary theoretical contributions:
(1) The dimensional invariant (oldest
surviving lineages sleep the least;
minimum sleep $=$ maximum survival horizon;
U-shaped survival--sleep relationship
with flat minimum).
(2) The two-strategy taxonomy
(Lock and Navigator).
(3) The unstable zone identification.
(4) The five sleep topologies as geometric
predictions of the landscape, not
empirical sleep types.
(5) The full three-channel sonar
(duration $+$ topology $+$ flexibility).
(6) Predictions P-INVARIANT-1 through
P-INVARIANT-7.\\
\textbf{Theory paper:}
\href{https://doi.org/10.5281/zenodo.19094935}
{\texttt{doi:10.5281/zenodo.19094935}}\\
\textbf{Repository:}
\href{https://github.com/Eric-Robert-Lawson/
attractor-oncology}
{\texttt{github.com/Eric-Robert-Lawson/
attractor-oncology}}\\
\textbf{ORCID:}
\href{https://orcid.org/0009-0002-0414-6544}
{0009-0002-0414-6544}\\
\textbf{Contact:}
\href{mailto:OrganismCore@proton.me}
{OrganismCore@proton.me}\\
\textbf{License:} CC BY 4.0\\
\textbf{Version:} 2.0 --- LOCKED AT
PRE-REGISTRATION.
}

\end{document}
